\documentclass[a4paper]{amsart}

\usepackage{color}
\usepackage[colorlinks=true, citecolor=red]{hyperref}  %needs to come before amsrefs package to work with \cite{}*{}
\usepackage{amssymb,latexsym}
\usepackage[alphabetic]{amsrefs}
\usepackage[mathscr]{eucal}
\usepackage{pdfsync}
\usepackage{graphicx}
\usepackage{epstopdf}
\DeclareGraphicsRule{.tif}{png}{.png}{`convert #1 `dirname #1`/`basename #1 .tif`.png}
\usepackage[all]{xy}
\usepackage{titletoc}
%\usepackage{txfonts}   %  \fint  for integral average

%\usepackage{showkeys}  %options notref  notcite 

%\input{rgb.tex}

%\textcolor{DarkRed}{text in DarkRed}


\CompileMatrices

\theoremstyle{plain}
\newtheorem{theorem}{Theorem}
\newtheorem{lemma}[theorem]{Lemma}
\newtheorem{proposition}[theorem]{Proposition}
\newtheorem{corollary}[theorem]{Corollary}

\theoremstyle{definition}
\newtheorem{definition}[theorem]{Definition}%[section]


\theoremstyle{remark}
\newtheorem{defn}[theorem]{Definition} 
\newtheorem{example}[theorem]{Example}
\newtheorem{remark}[theorem]{Remark}
\newtheorem{discussion}[theorem]{Discussion}

%\numberwithin{equation}{section}

%\newcommand{\abs}[1]{\lvert#1\rvert}

\DeclareMathOperator{\Lip}{Lip}
\DeclareMathOperator{\dist}{dist}
\DeclareMathOperator{\expected}{\mathbb{E}}
\DeclareMathOperator{\prb}{Prob}
\DeclareMathOperator{\spt}{spt}
\DeclareMathOperator{\lap}{\Delta}
\DeclareMathOperator{\Tr}{Tr}

\raggedbottom

\addtolength{\textwidth}{2cm}
\addtolength{\oddsidemargin}{-1cm}
\addtolength{\evensidemargin}{-1cm} 
\addtolength{\topmargin}{-1cm} 
\addtolength{\textheight}{2cm}

 
%%% BEGIN DOCUMENT
\begin{document}

\title{}  
\author{John E.\ Hutchinson}
\date{\today}


\maketitle


Suppose $ T:\mathbb{R}^2 \to \mathbb{R}^2 $ is linear.  Then either $ T $ has 2 (perhaps equal) real e-values or $ T $ has complex conjugate e-values.

The proof of the following proposition uses just linear algebra.

\begin{proposition} Suppose $ T:\mathbb{R}^2 \to \mathbb{R}^2 $ is linear with e-values $ \lambda e^{\pm i\theta} $ and $ \lambda > 0 $. Then $ \lambda ^{-1} T $ leaves some ellipse invariant.
\end{proposition}

\begin{proof} Using the same matrix as in the real case,  consider $ T $ as a (complex) linear map $ T: \mathbb{C}^2 \to \mathbb{C} ^2 $.
Diagonalising, there exists a linear $ S: \mathbb{C}^2 \to \mathbb{C} ^2 $ such that
\begin{equation} \label{mid}
S (\lambda ^{-1} T) S^{-1} = Q = \begin{bmatrix} e^{  i\theta}& 0 \\ 0 & e^{ - i\theta} \end{bmatrix}.
\end{equation} 
 $ Q $ fixes the unit sphere 
\begin{align*}
  \{ (z_1,z_2 ) \in \mathbb{C} ^2 : |z_1|^2 + |z_2|^2 = 1 \} &= \{ (x_1,y_1,x_2,y_2) : x_1^2 + y_1^2 + x_2^2 + y_2^2 = 1 \}  \\
   & =:S^3 \subset \mathbb{R} ^4.
\end{align*} 
  
Imbed $ \mathbb{R} ^ 2 \subset \mathbb{C} ^2 \cong \mathbb{R} ^ 4 $ via $ (x_1,x_2) \leftrightarrow (x_1,x_2) \leftrightarrow (x_1,0,x_2,0) $.  Let $ V = S( \mathbb{R}^2) \subset \mathbb{R} ^4 $.  Note $ V $ is a real 2-dimensional subspace and $ S $ is injective onto its image.

Let 
\[
 C = V \cap S^3  \subset V  \subset \mathbb{R} ^ 4,  \quad E = S^{-1}(C) \subset \mathbb{R}^2 .
 \]
Then $ C $ is a unit circle and $ E $ is an ellipse.

Then  
\[
Q(C) = C 
\]
since $  Q(C) $ is a circle from the third term in \eqref{mid}, and $ Q(C)\subset V $ from the first term in \eqref{mid}.  But there is only one unit circle in $ V $.

Applying $ S^{-1} $ to the left of the first two terms in \eqref{mid}, and applying this to $ C $,
\[ 
\lambda ^{-1}T S^{-1}(C) = S^{-1}Q(C) = S^{-1} (C).
\] 
That is,
\[
\lambda ^{-1}T (E) = E. \qedhere
\] 
\end{proof} 


 
\end{document} 

